\section{能源危机与解决路径}

\subsection{AI 正在引发能源危机}

Hassabis 坦率地承认了一个不太舒适的事实：\textbf{AI 正在引发能源危机}。大规模训练和推理所需的数据中心消耗着巨量的电力，全球数据中心的能源需求正在以前所未有的速度增长。

\begin{warningbox}{AI 的能源悖论}
AI 技术一方面在解决人类的各种问题，另一方面却在制造新的问题——对能源的巨大需求。如果不加以解决，AI 的发展可能会因为能源供给不足而减速，或者加剧气候变化。

这是一个典型的``你需要 AI 来解决的问题，正是 AI 制造的问题''式悖论。
\end{warningbox}

\subsection{AI 也能解决能源问题}

但硬币的另一面是：\textbf{AI 同时具备解决能源危机的潜力}。这包括：

\begin{itemize}
    \item \textbf{核聚变}：AI 可以加速等离子体控制和材料科学研究，推动可控核聚变的实现
    \item \textbf{电网优化}：AI 驱动的智能电网管理可以显著提高能源效率
    \item \textbf{新材料发现}：AI 可以帮助发现更高效的太阳能电池材料、储能材料等
    \item \textbf{能源管理}：数据中心本身的能源使用效率也可以通过 AI 优化
\end{itemize}

\begin{importantbox}{能源问题的辩证视角}
AI 引发了能源需求的激增，但 AI 同时也可能是解决能源问题的关键工具。这种``制造问题 $\to$ 解决问题''的循环，本质上是所有变革性技术的共同特征——蒸汽机既创造了工业污染，也最终推动了清洁能源技术的发展。

关键在于：\textbf{解决的速度必须快于问题恶化的速度}。如果 AI 能在 10 年内推动核聚变的突破，那么当前的能源消耗就是值得付出的``投资''。
\end{importantbox}

\subsection{本章小结}

AI 正在引发能源危机，但也具备解决这一危机的潜力。关键在于解决速度是否能快于问题恶化的速度。核聚变、新材料发现和电网优化是 AI 可以施力的方向。

%% ==================== Section 9 ====================
